\documentclass[11pt,a4paper]{article}
\usepackage[margin=27mm]{geometry}
\usepackage{fontspec,unicode-math,amsmath,amsthm,booktabs,longtable,hyperref}
\setmainfont{STIXGeneral}\setmathfont{STIX Two Math}\hypersetup{hidelinks}
\newtheorem{theorem}{Theorem}\newtheorem{lemma}[theorem]{Lemma}
\newtheorem{proposition}[theorem]{Proposition}
\newcommand{\PSL}{\operatorname{PSL}}\newcommand{\Aut}{\operatorname{Aut}}
\newcommand{\Cl}{\operatorname{Cl}}\newcommand{\Cen}{\operatorname{C}}
\title{Marked groups and complete locally projective quotient censuses}
\author{Shared proofs for the Hartley--Leemans original and erratum}\date{}
\begin{document}\maketitle
We recover the Hartley--Leemans results~\cite{source,erratum} by proving the marked universal group and local geometry for the dodecahedral--hemi-icosahedral polytope, then classify its quotients by a complete conjugacy test. The same argument supplies the original universal-polytope theorem and every correction in its later erratum. The first stage bounds the actual presented group, independently of its finite image. The second stage uses the established group and geometry to carry out the exhaustive subgroup calculation. This order removes a formerly retained universal-group premise from the erratum proof.

The complete table companion and the six whole required programs belong to this common theory. Both source-specific applications use them; neither relies on the other application. Classical quotient theory and named finite-group algorithms remain explicit external inputs. We do not claim a new principal theorem, historical priority or formal verification of those algorithms.

\section{A complete conjugacy test}
Let $G$ be finite, let $B\subseteq G$ be a forbidden set, and let $N\leq G$. Our convention is $x^g=g^{-1}xg$.
\begin{lemma}\label{avoid}
The following conditions are equivalent:
\[
 N^g\cap B=\varnothing\quad(g\in G),\qquad
 N\cap\bigcup_{b\in B}\Cl_G(b)=\varnothing.
\]
For the second condition it suffices to test one representative of every conjugacy class of elements of $N$. The admissible subgroups are closed under conjugation and passage to subgroups.
\end{lemma}
\begin{proof}
An element $b$ lies in $N^g\cap B$ exactly when $b=x^g$ for some $x\in N$ and $b\in B$. This is precisely membership of $x$ in a $G$-class meeting $B$. Conjugation within $N$ does not change its $G$-class, proving the representative test. Both closure assertions follow directly from the first condition.
\end{proof}

Suppose that $G$ is the marked group of a regular polytope, that $K$ is the stabilizer of a facet, and that $\omega$ is a central involution of $K$. Assume the quotient under consideration has flags indexed by the cosets of $N$, and that every intersection $N^g\cap K$ is either $1$ or $\langle\omega\rangle$. These assumptions concern the actual quotient and local sections, not just an arbitrary permutation action.
\begin{proposition}\label{facets}
Put $m=|N\cap\Cl_G(\omega)|$ and $c=|\Cen_G(\omega)|$. If $h$ and $d$ count the facet orbits with stabilizer of order two and one, respectively, then
\[
 h=\frac{2cm}{|K||N|},\qquad
 d=\frac{|G|}{|K||N|}-\frac h2.
\]
Moreover the full automorphism group of the quotient is $N_G(N)/N$.
\end{proposition}
\begin{proof}
Facets are cosets of $K$. A coset $Kg$ has an order-two stabilizer in $N$ exactly when $\omega^g\in N$. There are $cm$ choices of $g$ satisfying this condition: each of the $m$ possible conjugates has $c$ preimages. Centrality of $\omega$ in $K$ makes this condition independent of the representative of $Kg$, so there are $cm/|K|$ such facets. Their $N$-orbits have size $|N|/2$, which gives the formula for $h$. Counting all $|G|/|K|$ facets by their orbit sizes gives the formula for $d$.

On flags the marked generators act by adjacency. A type-preserving automorphism of the quotient commutes with these generators and hence with their generated transitive action. An equivariant bijection of $N\backslash G$ is determined by the image of $N$. That image is $Ng$ with the same stabilizer precisely when $g$ normalizes $N$. The resulting permutations are $N_G(N)/N$. Conversely, these permutations preserve every flag adjacency and therefore the quotient faces and incidence. This proves the automorphism assertion.
\end{proof}

\begin{lemma}\label{maxima}
Let a conjugacy-invariant collection of subgroups of a finite group be closed under subgroups. Form a directed graph on its subgroup conjugacy classes by recording, for each representative, the ambient conjugacy classes of all its maximal proper subgroups. A vertex is maximal in the collection exactly when it has no incoming edge.
\end{lemma}
\begin{proof}
An incoming edge exhibits a proper containing member. Conversely, if $N$ is properly contained in a member, choose a containing member minimal above $N$. Finiteness makes $N$ a maximal subgroup of that member. Conjugating the containing member to the chosen representative gives the required recorded edge. Repeated edges are counted only once.
\end{proof}


\section{Finite presentations and matching bounds}
\begin{lemma}\label{bound}
Let $W$ be a presented group, let $L\leq W$ have finite index $d$, and suppose that a presented group $T$ maps onto $L$. Suppose $T$ has a cyclic subgroup $C=\langle x\rangle$ of index $e$. If a complete integer relation matrix for $C/[C,C]$, on the specified generator $x$, has one column with nonzero gcd $m$, then
\[
 |W|\leq dem.
\]
If $W$ maps onto a finite marked group of order $dem$, that epimorphism is an isomorphism.
\end{lemma}
\begin{proof}
A group generated by one element is abelian. Thus $C=C/[C,C]\cong\mathbb Z/m\mathbb Z$: its defining relation subgroup in $\mathbb Z$ is generated by the matrix entries. Hence $|T|=em$, while surjectivity gives $|L|\leq em$. Coset partitioning now gives $|W|=d|L|\leq dem$. A surjection onto a group of this upper-bound order has trivial kernel.
\end{proof}
The hypothesis concerns a complete subgroup relation matrix, not a matrix from an arbitrary finite quotient. In our calculation it is computed by GAP's modified Todd--Coxeter algorithm for the subgroup generated by the actual element $x$ of the presented group. The simpler finite permutation image is used only on the opposite side of the inequality. GAP's exact presentation algorithms and ACE's completed coset enumeration are retained computational premises; neither a list of permutations satisfying the relators nor a closed table with no verified subgroup enumeration would replace them.

\begin{lemma}\label{prime}
Let $K$ be finite and let $B\subseteq K$. Suppose that for every prime $p\mid |K|$ and every conjugacy class of subgroups of order $p$, some member meets $B$. Then every nontrivial subgroup $N\leq K$ has a conjugate meeting $B$.
\end{lemma}
\begin{proof}
Cauchy's theorem gives a subgroup $C\leq N$ of prime order. By hypothesis $C^g\cap B\ne\varnothing$ for some $g$, and this intersection is contained in $N^g\cap B$.
\end{proof}
There is a complete finite test without enumerating every subgroup: form a cyclic subgroup from each element-class representative of order $p$, identify these subgroups up to conjugacy, and test their conjugates against the elements of $B$ of order $p$. Every subgroup of order $p$ occurs because any of its nonidentity elements generates it. This proof does not require all elements of order $p$ to be conjugate.


\section{The literal universal group}
Let $W$ have generators $s_0,s_1,s_2,s_3$ and relations
\begin{equation}\label{W}
 s_i^2=1,\quad(s_0s_1)^5=(s_1s_2)^3=(s_2s_3)^5=1,
 \quad[s_i,s_j]=1\ (|i-j|>1),\quad(s_1s_2s_3)^5=1.
\end{equation}
Put $\omega=(s_0s_1s_2)^5$ and $\nu=\omega s_3$.
\begin{theorem}\label{structure}
The literal presented group is marked-isomorphic to
\[
 W\cong J_1\times\PSL_2(19),\qquad |W|=600415200.
\]
The normal closures $N'=\langle\!\langle\nu^3\rangle\!\rangle$ and $N''=\langle\!\langle\omega\rangle\!\rangle$ are respectively the $\PSL_2(19)$ and $J_1$ factors. In particular $W/N'\cong J_1$, $W/N''\cong\PSL_2(19)$, and $\nu$ has order six.
\end{theorem}
\begin{proof}
Use exactly the source subgroup
\[
 L=\langle u_1,u_2,u_3\rangle,\quad
 u_1=s_2s_1,\quad u_2=(s_0s_1)^2s_2s_3s_0,\quad
 u_3=s_0s_1s_0s_3s_2s_3s_1.
\]
The complete ACE enumeration on presentation \eqref{W} and these three subgroup words has index $20$. Reduced Reidemeister--Schreier rewriting followed by tracked Tietze transformations produces two generators $X,Y$ and the exact relators
\begin{equation}\label{T}
 Y^3=1,\qquad XYXY^{-1}XYX^{-1}YX=1.
\end{equation}
The tracked images of the three primary subgroup generators are
\[
 u_1=X^{-1}Y^{-1}X,\quad u_2=Y^{-1}X^{-1}Y,
 \quad u_3=Y^{-1}X^{-1}Y^2X.
\]
Consequently the source words $x=u_2^{-1}u_3u_1$ and $y=u_1u_3^{-1}$ become $X$ and $X^{-1}Y^{-3}XY=Y$. This checks the marked source presentation, not merely an abstract presentation of some subgroup. Both the complete initial presentation and all tracked generator images are saved.

In the group $T$ defined by \eqref{T}, modified Todd--Coxeter enumeration for the specified subgroup $\langle X\rangle$ gives index $175560$. Its one-column abelianized subgroup relation matrix has $290$ nonzero rows, all divisible by $171$, and contains the row $-171$. Thus its gcd is exactly $171$. Lemma~\ref{bound} gives
\[
 |W|\leq20\cdot175560\cdot171=600415200.
\]
The complete code performs these calculations in the presented groups before introducing any finite image.

For the lower bound we reuse the complete Janko marking search, printed in full in the accompanying common program appendix. It constructs four elements $j_i$ in the named TomLib model of $J_1$. The other four markers $p_i$ are the literal permutations on twenty points printed on source page six. The exact finite checks show
\[
 \langle(j_i,p_i):0\leq i\leq3\rangle=J_1\times P,
 \quad |J_1|=175560,\quad P\cong\PSL_2(19),\quad |P|=3420.
\]
Every relation in \eqref{W} is checked in this product. The named $J_1$ model is retained library data; the twenty-point group is additionally compared by an actual isomorphism with GAP's $\PSL_2(19)$ construction. The resulting epimorphism therefore has the upper-bound order and is an isomorphism.

In the Janko coordinates the images of $\omega,\nu$ have orders $2,3$; in the projective coordinates they have orders $1,2$. Thus $\omega$ belongs to the Janko factor and $\nu^3$ belongs to the projective factor, both nontrivially. Simplicity of the factors, also checked in the finite models, makes their normal closures the entire indicated factors. This proves the marked quotient presentations, since adding the respective word is exactly quotienting by its normal closure. The two coordinate orders of $\nu$ give order six.
\end{proof}

The source's twenty-coset action is retained literally in the data and has kernel $N''$. Our direct-product construction supplies the faithful action on $266+20=286$ points described in the source; it is not claimed to reproduce the unavailable historical permutation file. The preceding proof also recovers the source's two-generator presentation of $L$, its order $30020760$, and the composition factors. It avoids any reliance on that file, \texttt{g.txt}.


\section{Geometry and the two distinguished quotients}
We use the classical marked Coxeter group $[5,3]$ of order $120$ for the dodecahedron, its reversed marking for the icosahedron, and the antipodal quotient of order $60$ for the corresponding hemi-polyhedron. The central antipode is the fifth power of the product of the three distinguished reflections. We also use the standard string-C-group/polytope correspondence and the general semisparse quotient correspondence. These are retained geometric inputs, not the target universal-group theorem.

\begin{theorem}\label{geometry}
The group in Theorem~\ref{structure}, with its displayed marking, is a string C-group. Its regular polytope $\mathcal P$ has dodecahedral facets and hemi-icosahedral vertex figures. It is the universal polytope with these local types. The quotient with group $J_1$ has $1463$ dodecahedral facets and $2926$ vertices; the quotient with group $\PSL_2(19)$ is the self-dual $57$-cell with hemi-dodecahedral facets and hemi-icosahedral vertex figures. Both distinguished quotients have no proper polytope quotient under the semisparse convention.
\end{theorem}
\begin{proof}
For each of the three actual markings, the required program constructs the subgroup on every subset of the four markers and checks all $16^2=256$ intersection equalities. The markers are nonidentity involutions, their adjacent products have exact orders $5,3,5$, and the nonadjacent pairs commute. These checks establish the entire string C-property, not only its two endpoint intersections.

The endpoint orders are respectively $(120,60)$ for $W$ and $J_1$, and $(60,60)$ for the projective quotient, where the first entry is the facet group. Their Coxeter presentations give surjections from the classical local groups just described, and equality of their orders makes them marked isomorphisms. Thus their local polytopes are the claimed ones; an abstract group isomorphism to $C_2\times A_5$ alone would not suffice. The string-C-group correspondence constructs the regular polytopes. Group order divided by the appropriate endpoint order gives $1463$, $2926$ and $57$. Reversing the four actual projective markers extends to a bijective group homomorphism, explicitly checked by the program, and hence proves the asserted self-duality.

For universality, take any polytope with the indicated local types and let its involutory flag adjacencies be $r_i$. Nonadjacent adjacencies commute; the polygon sections give the three adjacent relations. At each vertex the hemi-icosahedral flag action gives $(r_1r_2r_3)^5=1$. Thus \eqref{W} acts transitively on the flags, and the stabilizer of a chosen flag identifies the polytope with the corresponding quotient of $\mathcal P$ by the general quotient correspondence. This applies without assuming the target regular. It also applies when some facets have been replaced by their antipodal quotients, since their adjacencies still satisfy the same facet relations. Every section-regular locally projective polytope of type $\{5,3,5\}$ is of one of these local types or has dual of one of them. This proves the precise classification coverage claimed in the source; the mixed-facet quotients below are additional, non-section-regular members of the full quotient census.

For the final assertion retain the classical fact that the hemi-icosahedron has no proper quotient and Hartley's general endpoint semisparseness theorem. Applied to these local groups, it says that a subgroup $N$ of the Janko marked group is semisparse exactly when
\[
 N^g\cap H_0H_3\subseteq\{1,\omega\}\quad\hbox{for all }g;
\]
in the projective quotient the right side is $\{1\}$. The facet alternatives are respectively the dodecahedron and its antipodal quotient, or only the hemi-dodecahedron. For each prime and every conjugacy class of prime-order subgroups, the finite program finds a member meeting the corresponding forbidden endpoint product. Lemma~\ref{prime} therefore excludes every nontrivial semisparse subgroup.

This calculation also checks the source's explicit Janko witnesses: in its marker indices the words
\[
 s_1s_2,\quad s_0s_1,\quad s_3s_1s_2s_1s_0,\quad
 s_3s_2s_1s_0s_1s_0,\quad
 s_1s_3s_2s_1s_3s_0s_1s_0s_2s_1s_0
\]
belong to $H_0H_3$ and have orders $3,5,19,11,7$. Its unique involution class and its prime-order odd Sylow groups give the original short Cauchy argument as well. For the projective quotient the orders tested are $2,3,5,19$; the test is on prime-order subgroups, so does not incorrectly require a Sylow $3$-subgroup to have order three.
\end{proof}


\section{The complete quotient census}
Theorems~\ref{structure} and~\ref{geometry} establish the universal marked group and its local sections before any subgroup enumeration is used. Put
\[
 H_0=\langle s_1,s_2,s_3\rangle,\qquad H_3=\langle s_0,s_1,s_2\rangle.
\]
The retained general endpoint theorem and the local quotient identifications in Theorem~\ref{geometry} give
\begin{equation}\label{criterion}
 N\text{ is semisparse}\quad\Longleftrightarrow\quad
 N^w\cap H_0H_3\subseteq\{1,\omega\}\quad(w\in W).
\end{equation}
The general quotient correspondence identifies these subgroup conjugacy classes with the quotient polytopes, and their inclusions with coverings. These general geometric inputs remain external; the particular group and its local marked identifications have been proved above. No assertion that an arbitrary polytopal orbit poset implies semisparseness is used.

The required program reconstructs the same direct-product marking as in Theorem~\ref{structure}, using the full Janko marking search and the source's literal twenty-point permutations. It checks all defining relators and both endpoint orders. Apply Lemma~\ref{avoid} to
\[
 B=H_0H_3\setminus\{1,\omega\}.
\]
The product has $1200$ elements, meeting $47$ of the $180$ element conjugacy classes of $W$ after the two permitted elements are removed. Every representative of every element class in each subgroup is tested against this full forbidden-class set. For a surviving subgroup, \eqref{criterion} also guarantees exactly the local alternatives required by Proposition~\ref{facets}. All vertex figures remain hemi-icosahedra; its $d$ facets are dodecahedra and its $h$ facets hemi-dodecahedra.

\begin{theorem}\label{census}
The group $W$ has $137$ conjugacy classes of semisparse subgroups, of which $26$ are maximal. The corresponding quotient polytopes comprise $65$ with the original facets and vertex figures, one with only hemi-dodecahedral facets, and $71$ with mixed facets. Thus $66$ are section regular. The eight rows removed in \cite{erratum} are exactly
\[
 13,\ 70,\ 78,\ 90,\ 91,\ 95,\ 111,\ 127.
\]
\end{theorem}
\begin{proof}
The required program uses GAP's algorithm enumerating every subgroup conjugacy class, yielding $1262$ representatives. It does not read the authors' old subgroup file, a semisparse list, or a list of quotient polytopes. Lemma~\ref{avoid} filters this exhaustive list to $137$ classes. Proposition~\ref{facets} gives $65$ rows with $h=0$, one with $d=0$, and $71$ with $d,h>0$. The two homogeneous facet cases are precisely section regular here, since the vertex figure is fixed and all lower sections are those of the two stated regular facets. The program separately enumerates maximal subgroups of every surviving representative. The resulting graph on ambient subgroup classes has $349$ distinct edges and, by Lemma~\ref{maxima}, $26$ maximal vertices.

For the source labels, all thirty generating sets in the original Table 1 are evaluated as literal words in this same marking. Exactly the four labelled $13,70,95,127$ fail the test; the other twenty-six are exactly the computed maximal classes. The four remaining deletions are identified from the printed inclusion rows as follows:
\begin{center}\begin{tabular}{cll}\toprule
Row&Printed subgroup&Identification in the literal parent\\\midrule
78&$C_{45}$&the unique subgroup class of order $45$ in row 13\\
111&$C_{15}$&the unique subgroup class of order $15$ in row 13\\
90&$C_{30}$&the unique cyclic subgroup class of order $30$ in row 70\\
91&$C_3\times D_{10}$&the order-$30$ nonabelian child of row 70\\
&&with normalizer quotient of order $12$\\\bottomrule\end{tabular}\end{center}
Here $D_{10}$ has order ten. The other nonabelian order-$30$ child has normalizer quotient of order $36$ and is not silently identified with row 91. For each of the eight rows the output gives actual permutations $x,q,b$ with
\[
 x\in N,\quad q\in W,\quad x^q=b\in H_0H_3\setminus\{1,\omega\}.
\]
These are explicit failures of the required universal conjugate condition.

Finally, the $137$ surviving historical rows, using all their printed maximal-subgroup links and facet counts, are matched bijectively to the computed $137$-vertex, $349$-edge graph. The thirty literal generating sets anchor the maximal rows; every surviving row is below a surviving maximum. The exact graph isomorphism, orders and both facet counts are checked, rather than only the grand total. All $137$ normalizer quotients also agree with the printed automorphism columns, after elementary direct-product identities. These include $D_{12}\cong C_2\times S_3$, $D_{20}\cong C_2\times D_{10}$, splitting coprime cyclic factors, and $\operatorname{GL}_2(4)\cong C_3\times A_5$: in the last identity the scalar center has order three, meets $\operatorname{SL}_2(4)\cong A_5$ trivially, and together these subgroups have the full order $180$. The complete original row set minus these survivors is exactly the displayed eight labels. The required finite output and the group/quotient premises above therefore establish all assertions.
\end{proof}


\section{Complete tables, source interpretation and retained computation}
The companion \emph{The 137 surviving quotient classes} is an integral part of this common paper. Its two full tables contain every surviving historical label, subgroup order, both facet counts, all $349$ maximal-containment links, the actual computed subgroup descriptions and all normalizer quotients. It is included unchanged as \texttt{complete-census.pdf}; its five pages belong to the common acquisition, so neither application silently borrows another application's table. The exact subgroup generators and all eight forbidden-conjugate witnesses are supplied as data in addition to the readable tables. They are outputs of the exhaustive computation, not a supplied list on which the exhaustive algorithm depends.

\paragraph{What the source tables do and do not specify.}
The classification gives exact permutation subgroups, stronger data than semidirect-product abbreviations. The historical subgroup-description column is checked for order, while the full printed automorphism column is checked up to the elementary direct-product identities described in the census paper. We do not infer an abstract subgroup isomorphism merely from its abbreviated type and order. In particular source rows $38,39,40$ all use $19:9$-style notation, while the actual centers have orders $1,1,3$; the distinct action at row $40$ remains explicit. Exact representatives preserve all actions.

The accessible arXiv edition has these additional discrepancies, kept distinct from the erratum and from the unseen journal edition. Table~1's order-$70$ row labelled $58$ corresponds to row $63$ in the complete tables; their row $58$ instead has order $76$. Page~14 row $21$ prints $2^2\times(19:3)$, whereas its order is $456$ and its actual group is $2^3\times(19:3)$. Row $27$ prints $2\times19:3$, whereas its order is $342$ and its actual group is $2\times(19:9)$. Row $41$ has an elementary abelian normal subgroup of order eight and exponent two; its printed $8:(7:3)$ is not literally a cyclic-eight semidirect product under the source's stated notation. These exact finite distinctions, and the $38$--$40$ center distinction, are recorded separately. None is silently transferred to the journal text.

The source's numerical timing and former software history are preserved as historical reports, not mathematical results to reproduce. Its $145$, $30$, $69$, $70$ and $75$ census totals are replaced, respectively, by $137$, $26$, $65$, $66$ and $71$ in their stated categories. Its two named simple quotients, universal presentation, local types, covering scope and quotient invariants are all retained in the corrected principal contract.


\paragraph{Algorithms and their trust boundary.}
The complete common appendix contains the six distinct required files: the presentation calculation, marking construction, geometric tests, literal maximal-word evaluation, exhaustive subgroup census and Python table reconciliation. The shared marking file is printed once. The whole historical \texttt{source\_tables.json} is required comparison data. Neither the old unavailable \texttt{g.txt} nor a saved group workspace is used. GAP's exact group, subgroup, presentation and Tietze algorithms, ACE's completed coset enumeration, the named TomLib $J_1$ model and NetworkX's graph-isomorphism routine remain computational premises. Their actual versions and selected file hashes are retained; this does not claim a proof or a hash binding of the entire transitive installation. All complete code and data remain in the declared scope. There is no new Lean declaration or claim of end-to-end formal verification.

\paragraph{Attribution and proof organization.}
The principal universal group, geometry and source tables are due to Hartley and Leemans. The present edition reuses the already reviewed conjugacy tests, orbit counts, literal marking computation, finite-presentation upper bound and full corrected census. The new organization changes no finite calculation: the group and local geometry are established before the census, and both are now part of the shared proof. Thus the erratum application needs no unpriced theorem from the original-paper application. The two applications specify their distinct source contracts and draw their conclusions from this complete common paper, its table companion and its full required programs. The original fifteen-page preprint and two-page erratum remain distinct manuscripts.

\begin{thebibliography}{9}
\bibitem{source} M. I. Hartley and D. Leemans, \emph{Quotients of a Universal Locally Projective Polytope of type $\{5,3,5\}$}, arXiv:math/0307075v1, 5 July 2003, fifteen pages. Journal reference: Math. Z. \textbf{247} (2004), 663--674; that edition has not been collated here.
\bibitem{erratum} M. I. Hartley and D. Leemans, erratum, Math. Z. \textbf{253} (2006), 433--434, doi:10.1007/s00209-005-0872-z.
\bibitem{arp} P. McMullen and E. Schulte, \emph{Abstract Regular Polytopes}, Cambridge University Press, 2002. Retained classical local groups and group/polytope correspondence.
\bibitem{hartley} M. I. Hartley, \emph{Simpler Conditions for Semisparse Subgroups}, source reference [7]. Its general endpoint criterion is retained through the exact statement on source page three; no fresh complete primary reading is claimed.
\end{thebibliography}
\end{document}
